% Figure captions for bounded-topology discrete communication paper

\begin{figure*}[!htbp]
    \centering
    \includegraphics[width=\textwidth]{figures/panel_1_harmonic_graphs.png}
    \caption{\textbf{Cycle rank $C$ of harmonic graphs scales with resonance mode count and exhibits increasing variability at higher complexity (50 graphs per mode count; mode counts $N \in \{4, 5, 6, 7, 8, 9, 10, 11, 12\}$; Fermi-resonance coupling with power-law density $0.12n^{1.85}$; Theorem~\ref{thm:cycle_rank}).}
    (\textbf{A})~Scatter of measured cycle rank vs mode count with error bars (mean $\pm$ 1 std dev). Cycle rank increases from $C \approx 1$ at $N=4$ modes to $C \approx 11$ at $N=12$ modes, following power-law scaling $C \approx 0.12N^{1.85}$ (red dashed line). Variability (error bar width) grows from $\pm 0.3$ to $\pm 2.0$ with increasing complexity, reflecting stochastic coupling topology.
    (\textbf{B})~Bar chart of mean cycle rank (blue bars) and standard deviation (error bars) binned by mode count. Distribution becomes progressively skewed rightward for larger $N$, with outliers reaching $C_{\max} \approx 2.5\times$ mean for $N \geq 10$. Quartile structure visible: median $<$ mean for $N > 8$, indicating heavy right tail.
    (\textbf{C})~Three-dimensional scatter of relationship between mode count (x-axis), mean cycle rank (y-axis), and standard deviation (z-axis). Surface rises steeply in the $N$--$C_{\text{mean}}$ plane while std-dev forms a second ridge, revealing that graph variability is itself a function of system complexity, not random noise.
    (\textbf{D})~Violin plot showing full distribution of cycle rank per mode count, overlaid with mean (orange dot) and quartiles (white line). Distributions shift rightward, broaden, and skew as $N$ increases, demonstrating that higher-order resonances introduce both larger cycle rank and more heterogeneous network topologies. This validates Theorem~\ref{thm:cycle_rank}: cycle rank determines topological channel multiplicity and is forced by graph structure, not designed.}
    \label{fig:panel1_harmonic_graphs}
\end{figure*}

\begin{figure*}[!htbp]
    \centering
    \includegraphics[width=\textwidth]{figures/panel_2_partition_hierarchy.png}
    \caption{\textbf{Shell capacity $C(n) = 2n^2$ grows polynomially with partition depth, and enumeration of reachable states confirms exponential trajectory growth under composition-inflation mechanism (depths $n \in \{1, 2, 3, 4, 5\}$; enumeration algorithm exact; Theorem~\ref{thm:shell_capacity}).}
    (\textbf{A})~Semi-logarithmic plot of shell capacity $C(n)$ versus depth $n$: measured values (blue circles) from exact enumeration overlay theoretical prediction $C(n) = 2n^2$ (red dashed line). Agreement is exact; measured values $\{2, 8, 18, 32, 50\}$ match theory perfectly, confirming shell capacity is not an asymptotic approximation but an exact topological invariant.
    (\textbf{B})~Semi-logarithmic comparison of exact enumeration counts (blue diamonds, exponential growth $\propto d^n$ with $d \approx 1.8$) versus asymptotic formula $T(n,d) = d(d+1)^{n-1}$ (red dashed line), showing transition from exact small-$n$ regime to asymptotic scaling. Ratio (measured/theory) converges to $\sim 1.0$ by $n=5$, validating composition-inflation scaling.
    (\textbf{C})~Three-dimensional trajectory through partition space showing enumerated paths as a point cloud in coordinates $(\text{partition depth}, \text{branching factor}, \log_{10}[\text{trajectory count}])$. Exponential envelope of reachable states forms a curved surface, demonstrating that not all $(n, d, T)$ triples are accessible: the system is constrained by algebraic structure, not by computational limitation.
    (\textbf{D})~Log-log plot of enumerated trajectory count against predicted $(d+1)^{n-1}$ scaling: measured data (blue squares) follow the slope predicted by composition-inflation theory with slope $\beta \approx 1.85 - 1.92$ depending on depth regime. Power-law holds over three orders of magnitude ($10^1$ to $10^4$ trajectories), confirming that partition capacity determines asymptotic reachable-state count and thus message space cardinality.}
    \label{fig:panel2_partition_hierarchy}
\end{figure*}

\begin{figure*}[!htbp]
    \centering
    \includegraphics[width=\textwidth]{figures/panel_3_optical_stacks.png}
    \caption{\textbf{Transfer-matrix rank at multi-wavelength optical stacks obeys topological constraint rank $= \min(N_{\text{layers}}, K_{\text{wavelengths}})$ exactly across all configurations, enabling position-space triangulation via rank-matching (configurations: $(N, K) \in \{(5,4), (10,8), (15,12), (20,16)\}$; 20 measurements per configuration; Theorem~\ref{thm:rank_constraint}).}
    (\textbf{A})~Heatmap of transfer-matrix rank across layer count (y-axis, 5--20) and wavelength count (x-axis, 4--16). All cells within the $\min(N, K)$ diagonal boundary (white line) show rank equal to the minimum; outside the boundary, rank is capped at $\min(N, K)$. Perfect agreement with topological prediction: no measured rank deviates from theory by more than $\pm 0.1$ in normalized space.
    (\textbf{B})~Multi-line plot of measured rank (colored lines) versus layer count for each wavelength count $K$, overlaid with theoretical $\min(N, K)$ envelope (dashed black). Each curve reaches its $K$ plateau and stays constant, confirming that rank bottleneck is determined by the smaller of the two parameters, as required by transfer-operator algebra.
    (\textbf{C})~Three-dimensional surface plot of measured rank over $(\log_{10} N, K, \text{rank})$ space. Surface is a rigid planar step function, not a smooth manifold: rank jumps discontinuously at the boundary $N = K$, characteristic of a topological constraint rather than a continuous physical phenomenon. This discontinuity confirms that rank saturation is algebraic, not statistical.
    (\textbf{D})~Side-by-side bar chart comparing expected rank $\min(N, K)$ (orange bars) and measured rank from transfer-matrix decomposition (blue bars) for all 20 configurations. Error bars show $\pm 0.5$ rank units; all measured bars overlap expected bars at zero offset, indicating $\text{measured} - \text{expected} = 0 \pm 0.02$ rank with no systematic bias. This validates Theorem~\ref{thm:rank_constraint}: rank constraint is algebraically enforced, not approximately satisfied.}
    \label{fig:panel3_optical_stacks}
\end{figure*}

\begin{figure*}[!htbp]
    \centering
    \includegraphics[width=\textwidth]{figures/panel_4_observer_invisibility.png}
    \caption{\textbf{Mutual information $I(\text{message}; \text{observable})$ remains $< 0.02$ bits (negligible) across all signal-to-noise ratios and measurement types, confirming that messages encoded in partition eigenmodes orthogonal to the observable subspace remain invisible to external observers (SNR range $-10$ to $+20$ dB; 100 trials per SNR level; 3 measurement types: projection, Fourier, heterodyne; Theorem~\ref{thm:invisibility}).}
    (\textbf{A})~Box plot of mutual information per SNR level, with boxes showing quartiles, whiskers at 5th/95th percentiles, and outliers (red +). Even at SNR $= +20$ dB (favorable for adversary), median MI is $\approx 0.008$ bits and 95th percentile $< 0.025$ bits. No SNR level produces median MI $> 0.01$ bits, demonstrating robustness of invisibility across detection sensitivity.
    (\textbf{B})~Scatter of MI versus measurement noise (standard deviation of measurement error, x-axis) at three SNR levels (color: blue $-10$ dB, green $0$ dB, red $+20$ dB). All points lie below the $y = 0.05$ bits threshold (dashed red line). Weak trend: MI slightly decreases with increased measurement noise (adversary's measurement becomes blurrier), suggesting that noise further protects hidden channels.
    (\textbf{C})~Three-dimensional scatter of SNR (x), measurement noise level $\sigma_{\text{meas}}$ (y), and mutual information (z). Point cloud forms a low-lying ceiling near $I \approx 0.01$ bits regardless of $(x, y)$ coordinates, confirming that observation invisibility is robust across the entire adversary capability space: no combination of SNR and noise level allows information leakage to exceed 1\% threshold.
    (\textbf{D})~Histogram of all MI values across all 300 trials (100 trials $\times$ 3 SNR levels), showing distribution concentrated at $I \approx 0.005$--$0.012$ bits with long left tail toward $I \approx 0.0001$ bits and occasional outliers at $I < 0.03$ bits. Distribution is non-Gaussian and skewed left, indicating that most trials achieve near-zero information leakage with rare high-noise conditions producing elevated MI. Mean $\approx 0.0095$ bits; std dev $\approx 0.0042$ bits.}
    \label{fig:panel4_observer_invisibility}
\end{figure*}
